\documentclass[12pt]{article}

\usepackage[utf8]{inputenc}
\usepackage[T1]{fontenc}
\usepackage[margin=0.75in]{geometry}
\usepackage{array}
\usepackage{enumitem}
\pagestyle{empty}

\newcolumntype{C}{>{\centering\arraybackslash}p{1.15cm}}
\newcommand{\rowh}{\rule{0pt}{0.82cm}}

\usepackage[table]{xcolor}

\newcommand{\off}{\cellcolor{black!88}}

\begin{document}

\begin{center}
{\LARGE\bfseries Fill in the table}\\[4pt]
{\large Two strings of DNA letters. Where do they differ?}
\end{center}

\bigskip
\noindent
The reference is \textbf{GATTACA}. The read is \textbf{GATACA}. One letter is
missing from the read. The table finds out which one.

\subsection*{The rule for one cell}

Look at three neighbours and work out three numbers:

\begin{center}
\renewcommand{\arraystretch}{1.4}
\begin{tabular}{@{}lll@{}}
from the \textbf{upper-left} cell: & same letters: add 2 & different letters: subtract 1 \\
from the cell \textbf{above}: & subtract 2 & \\
from the cell to the \textbf{left}: & subtract 2 & \\
\end{tabular}
\end{center}

\noindent
Write the \textbf{biggest} of the three. If the biggest is below 0, write
\textbf{0}.

\medskip
\noindent
The border, in italics, is all zeros. Two cells are done for you:
\begin{itemize}[nosep]
\item G against G: upper-left is 0, same letters, so $0 + 2 = 2$. Above gives $-2$, left gives $-2$. Write 2.
\item G against A: upper-left gives $0 - 1 = -1$, above gives $-2$, left gives $2 - 2 = 0$. Write 0.
\end{itemize}

\medskip
\begin{center}
\renewcommand{\arraystretch}{1}
\begin{tabular}{|C|C||C|C|C|C|C|C|C|}
\hline
\rowh &  & \textbf{G} & \textbf{A} & \textbf{T} & \textbf{T} & \textbf{A} & \textbf{C} & \textbf{A} \\ \hline
\rowh & \textit{0} & \textit{0} & \textit{0} & \textit{0} & \textit{0} & \textit{0} & \textit{0} & \textit{0} \\ \hline\hline
\rowh \textbf{G} & \textit{0} & 2 & 0 & & & & & \\ \hline
\rowh \textbf{A} & \textit{0} & & & & & & & \\ \hline
\rowh \textbf{T} & \textit{0} & & & & & & & \\ \hline
\rowh \textbf{A} & \textit{0} & & & & & & & \\ \hline
\rowh \textbf{C} & \textit{0} & & & & & & & \\ \hline
\rowh \textbf{A} & \textit{0} & & & & & & & \\ \hline
\end{tabular}
\end{center}

\subsection*{When the table is full}

\begin{enumerate}[nosep]
\item What is the biggest number in the table, and where is it?
\item Start there. Ask which neighbour that number came from, and step to it.
      Repeat until you reach a 0. Draw the path.
\item One step of your path goes \textbf{left} instead of diagonally. That is
      the missing letter. Which one is it?
\item Which cells did you fill in that the path never came near?
\end{enumerate}

\newpage

\begin{center}
{\LARGE\bfseries Answer}
\end{center}

\bigskip
\begin{center}
\renewcommand{\arraystretch}{1}
\begin{tabular}{|C|C||C|C|C|C|C|C|C|}
\hline
\rowh &  & \textbf{G} & \textbf{A} & \textbf{T} & \textbf{T} & \textbf{A} & \textbf{C} & \textbf{A} \\ \hline
\rowh & \textit{0} & \textit{0} & \textit{0} & \textit{0} & \textit{0} & \textit{0} & \textit{0} & \textit{0} \\ \hline\hline
\rowh \textbf{G} & \textit{0} & \fbox{2} & 0 & 0 & 0 & 0 & 0 & 0 \\ \hline
\rowh \textbf{A} & \textit{0} & 0 & \fbox{4} & \fbox{2} & 0 & 2 & 0 & 2 \\ \hline
\rowh \textbf{T} & \textit{0} & 0 & 2 & 6 & \fbox{4} & 2 & 1 & 0 \\ \hline
\rowh \textbf{A} & \textit{0} & 0 & 2 & 4 & 5 & \fbox{6} & 4 & 3 \\ \hline
\rowh \textbf{C} & \textit{0} & 0 & 0 & 2 & 3 & 4 & \fbox{8} & 6 \\ \hline
\rowh \textbf{A} & \textit{0} & 0 & 2 & 0 & 1 & 5 & 6 & \fbox{10} \\ \hline
\end{tabular}
\end{center}

\bigskip
\begin{enumerate}
\item The biggest number is \textbf{10}, in the bottom-right cell.
\item The path is the boxed cells: $10 \leftarrow 8 \leftarrow 6 \leftarrow 4
      \leftarrow 2 \leftarrow 4 \leftarrow 2$.
\item The step from the boxed 4 in row T back to the boxed 2 in row A goes
      diagonally, and the step from that 2 back to the boxed 4 in row A goes
      \textbf{left}. So the read skips the first \textbf{T} of the reference:

\begin{center}
\ttfamily
\begin{tabular}{@{}l@{\ \ }l@{}}
reference: & G A T T A C A \\
read:      & G A \_ T A C A
\end{tabular}
\end{center}
\end{enumerate}

% --- attention sheet: generated by reference/make_fill_sheet.py ---
\newpage
\renewcommand{\rowh}{\rule{0pt}{0.7cm}}

\begin{center}
{\LARGE\bfseries Fill in the table, with attention}\\[4pt]
{\large Same rule. This time most of the table is switched off.}
\end{center}

\medskip
\noindent
The reference is \textbf{TATTATGCA}. The read is \textbf{TATTTGCA}. One letter is
missing from the read. We first slid the read along the reference to see where
the letters line up best, and kept only the cells near that place.
\textbf{Black cells are switched off: leave them empty, and count them as 0.}

\subsection*{The rule for one cell}

Look at three neighbours and work out three numbers:

\begin{center}
\renewcommand{\arraystretch}{1.4}
\begin{tabular}{@{}lll@{}}
from the \textbf{upper-left} cell: & same letters: add 2 & different letters: subtract 1 \\
from the cell \textbf{above}: & subtract 2 & \\
from the cell to the \textbf{left}: & subtract 2 & \\
\end{tabular}
\end{center}

\noindent
Write the \textbf{biggest} of the three. If the biggest is below 0, write
\textbf{0}.

\medskip
\noindent
The border, in italics, is all zeros. Two cells are done for you:
\begin{itemize}[nosep]
\item T against T: upper-left is 0, same letters, so $0 + 2 = 2$. Above gives $-2$, the cell to the left is black and counts as 0, so left gives $-2$. Write 2.
\item T against A: upper-left gives $0 - 1 = -1$. Above gives $-2$, left gives $2 - 2 = 0$. Write 0.
\end{itemize}

\medskip
\begin{center}
\renewcommand{\arraystretch}{1}
\begin{tabular}{|C|C||C|C|C|C|C|C|C|C|C|}
\hline
\rowh &  & \textbf{T} & \textbf{A} & \textbf{T} & \textbf{T} & \textbf{A} & \textbf{T} & \textbf{G} & \textbf{C} & \textbf{A} \\ \hline
\rowh & \textit{0} & \textit{0} & \textit{0} & \textit{0} & \textit{0} & \textit{0} & \textit{0} & \textit{0} & \textit{0} & \textit{0} \\ \hline\hline
\rowh \textbf{T} & \textit{0} & 2 & 0 &  & \off & \off & \off & \off & \off & \off \\ \hline
\rowh \textbf{A} & \textit{0} & \off &  &  &  & \off & \off & \off & \off & \off \\ \hline
\rowh \textbf{T} & \textit{0} & \off & \off &  &  &  & \off & \off & \off & \off \\ \hline
\rowh \textbf{T} & \textit{0} & \off & \off & \off &  &  &  & \off & \off & \off \\ \hline
\rowh \textbf{T} & \textit{0} & \off & \off & \off & \off &  &  &  & \off & \off \\ \hline
\rowh \textbf{G} & \textit{0} & \off & \off & \off & \off & \off &  &  &  & \off \\ \hline
\rowh \textbf{C} & \textit{0} & \off & \off & \off & \off & \off & \off &  &  &  \\ \hline
\rowh \textbf{A} & \textit{0} & \off & \off & \off & \off & \off & \off & \off &  &  \\ \hline
\end{tabular}
\end{center}

\subsection*{When the white cells are full}

\begin{enumerate}[nosep]
\item What is the biggest number in the table, and where is it?
\item Start there. Ask which neighbour that number came from, and step to it.
      Repeat until you reach a 0. Draw the path.
\item One step of your path goes \textbf{left} instead of diagonally. That is
      the missing letter. Which one is it?
\item How many cells did you fill in? How many would the whole table have?
\end{enumerate}

\newpage

\begin{center}
{\LARGE\bfseries Answer}
\end{center}

\bigskip
\begin{center}
\renewcommand{\arraystretch}{1}
\begin{tabular}{|C|C||C|C|C|C|C|C|C|C|C|}
\hline
\rowh &  & \textbf{T} & \textbf{A} & \textbf{T} & \textbf{T} & \textbf{A} & \textbf{T} & \textbf{G} & \textbf{C} & \textbf{A} \\ \hline
\rowh & \textit{0} & \textit{0} & \textit{0} & \textit{0} & \textit{0} & \textit{0} & \textit{0} & \textit{0} & \textit{0} & \textit{0} \\ \hline\hline
\rowh \textbf{T} & \textit{0} & \fbox{2} & 0 & 2 & \off & \off & \off & \off & \off & \off \\ \hline
\rowh \textbf{A} & \textit{0} & \off & \fbox{4} & 2 & 1 & \off & \off & \off & \off & \off \\ \hline
\rowh \textbf{T} & \textit{0} & \off & \off & \fbox{6} & 4 & 2 & \off & \off & \off & \off \\ \hline
\rowh \textbf{T} & \textit{0} & \off & \off & \off & \fbox{8} & \fbox{6} & 4 & \off & \off & \off \\ \hline
\rowh \textbf{T} & \textit{0} & \off & \off & \off & \off & 7 & \fbox{8} & 6 & \off & \off \\ \hline
\rowh \textbf{G} & \textit{0} & \off & \off & \off & \off & \off & 6 & \fbox{10} & 8 & \off \\ \hline
\rowh \textbf{C} & \textit{0} & \off & \off & \off & \off & \off & \off & 8 & \fbox{12} & 10 \\ \hline
\rowh \textbf{A} & \textit{0} & \off & \off & \off & \off & \off & \off & \off & 10 & \fbox{14} \\ \hline
\end{tabular}
\end{center}

\bigskip
\begin{enumerate}
\item The biggest number is \textbf{14}, in the bottom-right cell.
\item The path is the boxed cells: $14 \leftarrow 12 \leftarrow 10 \leftarrow 8 \leftarrow 6 \leftarrow 8 \leftarrow 6 \leftarrow 4 \leftarrow 2$.
\item The step from the boxed 8 in row T back to the boxed 6 in row T goes diagonally, and the step from that 6 back to the boxed 8 in row T goes \textbf{left}. So the read skips the second \textbf{A} of the reference:

\begin{center}
\ttfamily
\begin{tabular}{@{}l@{\ \ }l@{}}
reference: & T A T T A T G C A \\
read:      & T A T T \_ T G C A
\end{tabular}
\end{center}

\item 23 cells. The whole table has 72.
\end{enumerate}
% --- end of attention sheet ---

\end{document}
